% compile: xelatex AUN3_4095101.tex
\documentclass{mko3}
\begin{document}
\mkothreetitle

% ----- หมวดที่ 1 ข้อมูลทั่วไป (รูปแบบ มคอ.3 AUN-QA) -----
\coursecode{4095101}
\coursename{ปัญญาประดิษฐ์และการเรียนรู้ของเครื่อง}
\program{หลักสูตรวิศวกรรมศาสตรมหาบัณฑิต สาขาวิชาวิศวกรรมคอมพิวเตอร์และปัญญาประดิษฐ์}
\coursecredits{3 (2-2-5)}
\coursegroup{วิชาเฉพาะด้านบังคับ}
\yearlevel{นักศึกษาปริญญาโท — ภาคต้น (ตามแผนการศึกษา)}
\semester{1 / 2568}
\instructor{ผศ.ดร.พิศิษฐ์ นาคใจ และ ผศ.ดร.พัชรี มณีรัตน์}
\contact{pisit.nakjai@uru.ac.th, patcharee.manee@uru.ac.th}
\coursedesc{พื้นฐานของปัญญาประดิษฐ์และการเรียนรู้ของเครื่อง การแก้ปัญหาด้วยการค้นหา การจัดการและการวิเคราะห์ข้อมูล การเรียนรู้แบบมีผู้สอนและไม่มีผู้สอน เทคนิคการหาค่าเหมาะสมและการลดค่าความชันเชิงเกรเดียนต์ โครงข่ายประสาทเทียมและการเรียนรู้เชิงลึก การวัดประสิทธิภาพและปรับปรุงแบบจำลอง การประยุกต์ใช้งานจริงภายใต้หลักจริยธรรมและความรับผิดชอบต่อสังคม}
\courseinfotable
\begin{clodomaintable}
  \cloitem{CLO1}{อธิบาย หลักการพื้นฐานของปัญญาประดิษฐ์และการเรียนรู้ของเครื่อง ทั้งแบบมีผู้สอนและไม่มีผู้สอนได้อย่างถูกต้อง}{Cognitive}{2}
  \cloitem{CLO2}{ประยุกต์ใช้ เทคนิคการแก้ปัญหาด้วยการค้นหาและการวิเคราะห์ข้อมูล ในการพัฒนาแบบจำลองการเรียนรู้ของเครื่องได้}{Cognitive}{3}
  \cloitem{CLO3}{วิเคราะห์ ประสิทธิภาพของแบบจำลองโครงข่ายประสาทเทียมและการเรียนรู้เชิงลึก เพื่อเลือกใช้และปรับปรุงให้เหมาะสมกับปัญหาที่กำหนดได้}{Cognitive}{4}
  \cloitem{CLO4}{แสดงออก ถึงความตระหนักในจริยธรรมและความรับผิดชอบต่อสังคม ในการพัฒนาและประยุกต์ใช้ระบบปัญญาประดิษฐ์}{Affective}{3}
  \cloitem{CLO5}{พัฒนา ระบบปัญญาประดิษฐ์ โดยใช้เทคนิคการเรียนรู้ของเครื่องที่เหมาะสมในการแก้ปัญหาจริงได้อย่างมีประสิทธิภาพ}{Psychomotor}{3}
\end{clodomaintable}

\begin{cloplotable}{5}{ & \textbf{PLO1} & \textbf{PLO2} & \textbf{PLO3} & \textbf{PLO4} & \textbf{PLO5}}
  \textbf{CLO1} & X &   &   &   &   \\ \hline
  \textbf{CLO2} & X &   &   &   &   \\ \hline
  \textbf{CLO3} &   &   & X &   &   \\ \hline
  \textbf{CLO4} &   &   & X &   &   \\ \hline
  \textbf{CLO5} & X &   &   &   &   \\ \hline
\end{cloplotable}

\begin{bloomtable}
  \bloomrow{CLO1}{2}{}{}{}
  \bloomrow{CLO2}{3}{}{}{}
  \bloomrow{CLO3}{4}{}{}{}
  \bloomrow{CLO4}{}{3}{}{}
  \bloomrow{CLO5}{}{}{3}{3}
\end{bloomtable}

\begin{contentanalysistable}
  \contentrow{หลักการพื้นฐานของปัญญาประดิษฐ์และการเรียนรู้ของเครื่อง}{Cognitive 2}{อธิบาย}{ทั้งแบบมีผู้สอนและไม่มีผู้สอนได้อย่างถูกต้อง (CLO1) [PLO1]}{4}
  \contentrow{เทคนิคการแก้ปัญหาด้วยการค้นหาและการวิเคราะห์ข้อมูล}{Cognitive 3}{ประยุกต์ใช้}{ในการพัฒนาแบบจำลองการเรียนรู้ของเครื่องได้ (CLO2) [PLO1]}{4}
  \contentrow{ประสิทธิภาพของแบบจำลองโครงข่ายประสาทเทียมและการเรียนรู้เชิงลึก}{Cognitive 4}{วิเคราะห์}{เพื่อเลือกใช้และปรับปรุงให้เหมาะสมกับปัญหาที่กำหนดได้ (CLO3) [PLO3]}{4}
  \contentrow{จริยธรรมและความรับผิดชอบต่อสังคม}{Affective 3}{แสดงออก}{ในการพัฒนาและประยุกต์ใช้ระบบปัญญาประดิษฐ์ (CLO4) [PLO3]}{4}
  \contentrow{ระบบปัญญาประดิษฐ์}{Psychomotor 3}{พัฒนา}{โดยใช้เทคนิคการเรียนรู้ของเครื่องที่เหมาะสมในการแก้ปัญหาจริงได้อย่างมีประสิทธิภาพ (CLO5) [PLO1]}{4}
\end{contentanalysistable}

\mkocoursedesc

\begin{courseactivities}
  \actitem{จดบันทึก ตั้งคำถาม อภิปรายกลุ่ม และทำแบบฝึกหัด (สอดคล้อง CLO1)}
  \actitem{ฝึกปฏิบัติการวิเคราะห์ข้อมูล ทำโครงงานย่อย และนำเสนอผล (CLO2)}
  \actitem{พัฒนาโมเดล Deep Learning ทดสอบประสิทธิภาพ และปรับปรุงโมเดล (CLO3)}
  \actitem{อภิปรายกลุ่ม เขียนบทความสะท้อนคิด และนำเสนอมุมมองด้านจริยธรรม AI (CLO4)}
  \actitem{พัฒนาระบบ AI ทดสอบระบบ นำเสนอผลงาน และแลกเปลี่ยนกับผู้เชี่ยวชาญ (CLO5; IEL)}
\end{courseactivities}

\begin{teachingtable}
  \teachrow{CLO1}{บรรยายทฤษฎีและหลักการ - ยกตัวอย่างกรณีศึกษา - สาธิตการทำงานของ AI [IEL: ประสบการณ์เชิงบูรณาการ ลงมือปฏิบัติ พัฒนาท้องถิ่น]}{การสอบข้อเขียน - การนำเสนอหน้าชั้น - การมีส่วนร่วมในชั้นเรียน}
  \teachrow{CLO2}{สาธิตการใช้เครื่องมือ - แนะนำเทคนิคการวิเคราะห์ - ให้คำปรึกษาระหว่างปฏิบัติ [IEL: ประสบการณ์เชิงบูรณาการ ลงมือปฏิบัติ พัฒนาท้องถิ่น]}{การสอบปฏิบัติ - รายงานผลการวิเคราะห์ - การนำเสนอโครงงาน}
  \teachrow{CLO3}{สาธิตการพัฒนาโมเดล - วิเคราะห์กรณีศึกษา - ให้คำแนะนำการปรับปรุงโมเดล [IEL: ประสบการณ์เชิงบูรณาการ ลงมือปฏิบัติ พัฒนาท้องถิ่น]}{โครงงาน DL - รายงานผลการทดสอบ - การนำเสนอผลงาน}
  \teachrow{CLO4}{อภิปรายประเด็นจริยธรรม - วิเคราะห์กรณีศึกษา [IEL: ประสบการณ์เชิงบูรณาการ ลงมือปฏิบัติ พัฒนาท้องถิ่น]}{การมีส่วนร่วมในการอภิปราย - บทความสะท้อนคิด}
  \teachrow{CLO5}{ให้โจทย์ปัญหาจริง - ให้คำปรึกษาโครงงาน - เชิญผู้เชี่ยวชาญ [IEL: ประสบการณ์เชิงบูรณาการ ลงมือปฏิบัติ พัฒนาท้องถิ่น]}{โครงงานปลายภาค - การสาธิตระบบ - รายงานโครงงาน}
\end{teachingtable}

\begin{assessmenttable}
  \assessrow{สอบก่อนเรียน (ถ้ามี)}{---}{0}
  \assessrow{สอบระหว่างเรียน}{CLO1, CLO2, CLO3}{30}
  \assessrow{สอบปลายภาค}{CLO1, CLO2, CLO3, CLO4, CLO5}{40}
  \assessrow{รายงาน / โครงงาน / กิจกรรม}{CLO3, CLO4, CLO5}{30}
\end{assessmenttable}

\begin{weeklyplan}
  \weekrow{1}{บทนำ AI/ML และขอบเขตรายวิชา}{บรรยาย + ปฐมนิเทศ + แบบทดสอบก่อนเรียน (ถ้ามี)}{CLO1}{สอบก่อนเรียน}
  \weekrow{2}{การแก้ปัญหาด้วยการค้นหา (Search)}{บรรยาย + สาธิต + แบบฝึกหัด}{CLO1, CLO2}{แบบฝึกหัด}
  \weekrow{3}{การเรียนรู้แบบมีผู้สอน (Supervised Learning)}{บรรยาย + Lab สาธิต}{CLO1, CLO2}{แบบฝึกหัด}
  \weekrow{4}{การเรียนรู้แบบไม่มีผู้สอน (Unsupervised Learning)}{บรรยาย + Lab}{CLO1, CLO2}{รายงานย่อย}
  \weekrow{5}{การจัดการและวิเคราะห์ข้อมูล}{Workshop + ปฏิบัติการ}{CLO2}{สอบกลางภาค 1}
  \weekrow{6}{การหาค่าเหมาะสมและ Gradient Descent}{บรรยาย + Lab}{CLO2}{แบบฝึกหัด}
  \weekrow{7}{โครงข่ายประสาทเทียมพื้นฐาน}{บรรยาย + สาธิต}{CLO3}{แบบฝึกหัด}
  \weekrow{8}{การเรียนรู้เชิงลึก (Deep Learning)}{บรรยาย + Lab GPU}{CLO3}{รายงานกลางภาค}
  \weekrow{9}{การวัดและปรับปรุงประสิทธิภาพโมเดล}{บรรยาย + วิเคราะห์กรณีศึกษา}{CLO3}{สอบกลางภาค 2}
  \weekrow{10}{จริยธรรมและความรับผิดชอบต่อสังคมใน AI}{อภิปราย + กรณีศึกษา}{CLO4}{บทความสะท้อนคิด}
  \weekrow{11}{ออกแบบโครงงานและกำหนดโจทย์}{PBL + ให้คำปรึกษา}{CLO5}{ข้อเสนอโครงงาน}
  \weekrow{12}{พัฒนาโครงงาน (ช่วงที่ 1)}{Workshop}{CLO5}{รายงานความก้าวหน้า}
  \weekrow{13}{พัฒนาโครงงาน (ช่วงที่ 2)}{Workshop}{CLO5}{---}
  \weekrow{14}{ทดสอบและปรับปรุงระบบ}{สาธิต + ทดสอบ}{CLO5}{สาธิตระบบ}
  \weekrow{15}{นำเสนอโครงงานและสรุปรายวิชา}{นำเสนอ + สอบปลายภาค}{CLO1--CLO5}{สอบปลายภาค + โครงงาน}
\end{weeklyplan}

% ----- หมวดที่ 5 Rubric -----
\begin{rubricsection}
\begin{subrubric}{CLO1}{อธิบาย หลักการพื้นฐานของปัญญาประดิษฐ์และการเรียนรู้ของเครื่อง ทั้งแบบมีผู้สอนและไม่มีผู้สอนได้อย่างถูกต้อง}
  \rubricrow{4 (ดีเยี่ยม)}{อธิบายหลักการได้อย่างถูกต้องลึกซึ้ง เชื่อมโยงทฤษฎีกับบริบทวิชาชีพและวรรณกรรมที่เกี่ยวข้องได้อย่างเหมาะสม}
  \rubricrow{3 (ดี)}{อธิบายหลักการได้ถูกต้องส่วนใหญ่ ครอบคลุมประเด็นสำคัญตามมาตรฐานหลักสูตรปริญญาโท}
  \rubricrow{2 (พอใช้)}{อธิบายได้บางส่วน ยังขาดความลึกหรือการเชื่อมโยงเชิงวิชาการ}
  \rubricrow{1 (ต้องปรับปรุง)}{ไม่สามารถอธิบายได้อย่างถูกต้องหรือไม่ครบถ้วน}
\end{subrubric}

\begin{subrubric}{CLO2}{ประยุกต์ใช้ เทคนิคการแก้ปัญหาด้วยการค้นหาและการวิเคราะห์ข้อมูล ในการพัฒนาแบบจำลองการเรียนรู้ของเครื่องได้}
  \rubricrow{4 (ดีเยี่ยม)}{ประยุกต์ใช้ได้อย่างถูกต้องและเหมาะสม แก้ปัญหาซับซ้อนได้ด้วยเหตุผลและหลักฐานสนับสนุน}
  \rubricrow{3 (ดี)}{ประยุกต์ใช้ได้ถูกต้องส่วนใหญ่ ตรงตามวัตถุประสงค์ของงาน}
  \rubricrow{2 (พอใช้)}{ประยุกต์ใช้ได้บางส่วน ยังไม่เหมาะสมหรือไม่สมบูรณ์}
  \rubricrow{1 (ต้องปรับปรุง)}{ไม่สามารถประยุกต์ใช้ได้หรือใช้ผิดแนวทาง}
\end{subrubric}

\begin{subrubric}{CLO3}{วิเคราะห์ ประสิทธิภาพของแบบจำลองโครงข่ายประสาทเทียมและการเรียนรู้เชิงลึก เพื่อเลือกใช้และปรับปรุงให้เหมาะสมกับปัญหาที่กำหนดได้}
  \rubricrow{4 (ดีเยี่ยม)}{วิเคราะห์เชิงวิพากษ์อย่างเป็นระบบ ใช้หลักฐาน/วรรณกรรมสนับสนุน สรุปและเสนอแนวทางได้ชัดเจน}
  \rubricrow{3 (ดี)}{วิเคราะห์ได้ถูกต้องส่วนใหญ่ มีเหตุผลประกอบเพียงพอตามมาตรฐานปริญญาโท}
  \rubricrow{2 (พอใช้)}{วิเคราะห์ได้ตื้น หรือยังขาดเหตุผล/หลักฐานสนับสนุน}
  \rubricrow{1 (ต้องปรับปรุง)}{วิเคราะห์ไม่ถูกต้องหรือไม่สามารถวิเคราะห์ได้}
\end{subrubric}

\begin{subrubric}{CLO4}{แสดงออก ถึงความตระหนักในจริยธรรมและความรับผิดชอบต่อสังคม ในการพัฒนาและประยุกต์ใช้ระบบปัญญาประดิษฐ์}
  \rubricrow{4 (ดีเยี่ยม)}{แสดงทัศนคติ/จริยธรรมวิชาชีพได้ชัดเจน สม่ำเสมอ และสะท้อนคิดอย่างลึกซึ้ง}
  \rubricrow{3 (ดี)}{แสดงออกได้เหมาะสมส่วนใหญ่ มีการสะท้อนคิด}
  \rubricrow{2 (พอใช้)}{แสดงออกได้บางครั้ง ยังไม่ชัดเจนหรือไม่สม่ำเสมอ}
  \rubricrow{1 (ต้องปรับปรุง)}{ไม่แสดงออกหรือแสดงออกไม่เหมาะสม}
\end{subrubric}

\begin{subrubric}{CLO5}{พัฒนา ระบบปัญญาประดิษฐ์ โดยใช้เทคนิคการเรียนรู้ของเครื่องที่เหมาะสมในการแก้ปัญหาจริงได้อย่างมีประสิทธิภาพ}
  \rubricrow{4 (ดีเยี่ยม)}{ออกแบบ/พัฒนางานวิจัยหรือระบบได้สมบูรณ์ มีคุณภาพระดับบัณฑิตศึกษา นำไปใช้หรือเผยแพร่เชิงวิชาการได้}
  \rubricrow{3 (ดี)}{พัฒนาได้ดี ตรงวัตถุประสงค์และมาตรฐานที่ยอมรับได้}
  \rubricrow{2 (พอใช้)}{พัฒนาได้บางส่วน ยังไม่สมบูรณ์}
  \rubricrow{1 (ต้องปรับปรุง)}{ไม่สามารถพัฒนางานได้ตามที่กำหนด}
\end{subrubric}
\end{rubricsection}


\begin{learningmaterials}
  \matitem{เอกสารประกอบการบรรยาย (Slides) และคู่มือปฏิบัติการ Lab}
  \matitem{ตำราหลัก: Russell \& Norvig, \textit{Artificial Intelligence: A Modern Approach}; Goodfellow et al., \textit{Deep Learning}}
  \matitem{ชุดข้อมูลฝึกปฏิบัติ, Jupyter Notebook, Python (scikit-learn, TensorFlow/PyTorch)}
  \matitem{บทความวิชาการและกรณีศึกษาจาก IEEE/ACM Digital Library}
  \matitem{วิดีโอบรรยายย้อนหลังและสื่อออนไลน์ผ่านระบบ LMS มหาวิทยาลัย}
\end{learningmaterials}

\begin{supporttable}
  \supportrow{ห้องเรียน Smart Classroom}{1 ห้อง}{พร้อมใช้}{ทันสมัย}
  \supportrow{ห้องปฏิบัติการคอมพิวเตอร์}{1 ห้อง}{พร้อมใช้}{ทันสมัย}
  \supportrow{GPU Server / Cloud Computing Resources}{ครบถ้วน}{พร้อมใช้}{ทันสมัย}
  \supportrow{เอกสารประกอบการสอน, ตำรา, วารสารวิชาการ}{ครบถ้วน}{พร้อมใช้}{ทันสมัย}
\end{supporttable}

\signatureblock{ผศ.ดร.พิศิษฐ์ นาคใจ และ ผศ.ดร.พัชรี มณีรัตน์}{ผศ.ดร.พิทักษ์ คล้ายชม}

\end{document}